%%% In this section, you will describe all of the various artifacts that you will generate and maintain during the project life cycle. Describe the purpose of each item below, how the content will be generated, where it will be stored, how often it will be updated, etc. Replace the default text for each section with your own description. Reword this paragraph as appropriate.

\subsection{Major Documentation Deliverables}
\subsubsection{Project Charter}
* The project charter will be created during the first sprint period, and this will be our guiding contract between the team, the stakeholders, and the faculty sponsor(Professor Geiser). The first version will be delivered by the end of the first sprint. We will update if there is a significant change in the project's direction/goal. The final version will be delivered on the last sprint cycle which is in late April to early May of 2026

\subsubsection{System Requirements Specification}
* This will outline all the functional/non-functional requirements for the PSIEM system. The document will be updated each time new requirements are classified or old ones are modified. The final version should be submitted through Senior Design.

\subsubsection{Architectural Design Specification}
* This will describe the architecture of PSIEM, we will also include the choices of frameworks, cloud platform, and communication, and make sure they are up to date. The initial version should be completed by Sprint 3, and updates will occur until all the design requirements are agreed upon. The final Architectural Design Specification will be delivered near the end of Senior Design 1.

\subsubsection{Detailed Design Specification}
* This specification will go into more detail for the architecture, into module-level details, that will include algorithms we might use or specific interface designs that might be used. The initial draft will be prepared by Sprint 4 after the architectural design. The final Detailed Design should be given by Senior Design 2, when actual implementation begins to start.

\subsection{Recurring Sprint Items}
\subsubsection{Product Backlog}
Items will be organize and added to the product backlog by the specific sprint cycle by the current development cycle on the system overview (with full sub detail categories) diagram. We shall initialize group votes for group consensus on items that need to go on the product backlog as the project's components are split and broken down upon each team memmber's strongest area of expertise. 

Software that would be use to share and maintain the product backlog will be discord ( for Dev team ), google slides ( for presentational sprint cycles/review ), and lucid chart for team discussion and modifications. 

\subsubsection{Sprint Planning}
The sprint plan/cycles will be planned via team meetings/discussion at least 2 days prior to presentation period via Discord. For the first half senio design there is expected to be four Spint plan cycles and the later half there will be about 5 sprint plan cycles.

\subsubsection{Sprint Goal}
The scrum master ultimately describes the sprint goal but final draft of the sprint goal for the specific sprint will be finalized by a group consesus and agreement. The customer is the overall important asset in this project so the sprint goal will reflect the components we intend to implement in such that the customer will have access and understanding.

\subsubsection{Sprint Backlog}
Who decides which product backlog items make their way into the sprint backlog? How will the backlog be maintained (collaboration software, a "scrum board", etc.)?

\subsubsection{Task Breakdown}
Individual task will be assigned via each team member's confidence and overall experience or expertise in the suggest area of the product back items. Members can choose to also voluntarily help other members with their current task if need help. Documentation on task hours for each team memmber will be documented on an excel sheet with corresponding sprint cycles days.

\subsubsection{Sprint Burn Down Charts}
Who will be responsible for generating the burn down charts for each sprint? How will they be able to access the total amount of effort expended by each individual team member? What format will the burn down chart use (include an example burn down chart below).

\begin{figure}[h!]
    \centering
    \includegraphics[width=0.5\textwidth]{images/test_image}
    \caption{Example sprint burn down chart}
\end{figure}

\subsubsection{Sprint Retrospective}
How will the sprint retrospective be handled as a team? When will this discussion happen after each sprint? What will be documented as a group and as individuals, and when will it be due?

\subsubsection{Individual Status Reports}
What sort of status will be reported by each individual member, and how often will it be reported? What key items will be contained in the report?

\subsubsection{Engineering Notebooks}
How often will the engineering notebook be updated, at a minimum, by each team member? What is the minimum amount of pages that will be completed for each interval, and how long will that interval be? How will the team keep each member accountable? Who will sign of as a "witness" for each ENB page?

\subsection{Closeout Materials}
The following materials, in addition to major documentation deliverables, will be provided to the customer upon project closeout. Remove this paragraph from your draft, but leave the heading.

\subsubsection{System Prototype}
What will be included in the final system prototype? How and when will this be demonstrated? Will there be a Prototype Acceptance Test (PAT) with your customer? Will anything be demonstrated off-site? If so, will there be a Field Acceptance Test (FAT)?

\subsubsection{Project Poster}
What will be included on the poster, what will be the final dimensions, and when will it be delivered?

\subsubsection{Web Page}
What will be included on the project web page? Will it be accessible to the public? When will this be delivered? Will it be updated throughout the project, or just provided at closeout (at a minimum, you need to provide a simple web page at the end).

\subsubsection{Demo Video}
What will be shown in the demo video(s)? Will you include a B-reel footage for future video cuts? Approximately how long will the video(s) be, and what topics will be covered?

\subsubsection{Source Code}
How will your source code be maintained? What version control system will you adopt? Will source code be provided to the customer, or binaries only? If source code is provided, how will it be turned over to the customer? Will the project be open sourced to the general public? If so, what are the license terms (GNU, GPL, MIT, etc.). Where will the license terms be listed (in each source file, in a single readme file, etc.).

\subsubsection{Source Code Documentation}
What documentation standards will be employed? Will you use tools to generate the documentation (Doxygen, Javadocs, etc.). In what format will the final documentation be provided (PDF, browsable HTML, etc.)?

\subsubsection{Hardware Schematics}
Will you be creating printed circuit boards (PCBs) or wiring components together? If so, list each applicable schematic and what sort of data it will contain (PCB layout, wiring diagram, etc.). If your project is purely software, omit this section.

\subsubsection{CAD files}
Will the project involve any mechanical design, such as 3D printed or laser-cut parts? If so, what software will you use to generate the files and what file formats will you provide in your closeout materials (STL, STEP, OBJ, etc.). If your project is purely software, omit this section.

\subsubsection{Installation Scripts}
How will the customer deploy software to new installations? Will you provide installation scripts, install programs, or any other tools to improve the process? Will there be multiple scripts provided (perhaps separate scripts for the graphical front end and back end server software)? 

\subsubsection{User Manual}
Will you customer need a printed or digital user manual? Will they need a setup video? Decide now what will be provided and discuss.
